\chapter{\texorpdfstring{Derived Mealy Morphisms and Homotopical $1$--Cocycles}{Derived Mealy Morphisms and Homotopical 1-Cocycles}}
\label{ch:derived-mealy-morphisms}

\section{Orientation: from strict machines to derived executions}
\label{sec:derived-mealy-orientation}

In the previous chapter, Mealy morphisms were introduced as monad
morphisms in the bicategory of spans. Equivalently, a strict Mealy
morphism
\[
(P,\delta,\omega):\mathbb A\rightsquigarrow \mathbb B
\]
is a typed stateful transducer. Its states form a span
\[
A_0 \xleftarrow{p} P \xrightarrow{q} B_0,
\]
an input arrow
\[
a:u\to v
\]
of \(\mathbb A\) is admissible at a state \(x\in P\) precisely when
\[
p(x)=v,
\]
and the machine produces
\[
x
\xrightarrow{\;a/\omega(a,x)\;}
\delta(a,x),
\]
where
\[
p(\delta(a,x))=u,
\qquad
\omega(a,x):q(\delta(a,x))\to q(x)
\]
is an arrow of \(\mathbb B\).

Thus the input is processed contravariantly: arrows of \(\mathbb A\) are
lifted target-to-source along the state projection \(p:P\to A_0\). The
output component \(\omega\) is then a nonabelian categorical cocycle
recording how output interfaces change along these lifted executions.

The purpose of this chapter is to pass from strict span-theoretic
representatives to a homotopy-invariant calculus of derived
correspondences. The guiding principle is:

\begin{quote}
A derived Mealy morphism is not merely a strict Mealy machine modulo
weak equivalence. It is a homotopy-coherent transducer: its rigid state
object is replaced by a derived execution object, and its strict output
law is replaced by a derived nonabelian output \(1\)-cocycle.
\end{quote}

The strict picture
\[
A_0 \xleftarrow{p} P \xrightarrow{q} B_0
\]
will be replaced by a derived correspondence
\[
\mathbb A
\longleftarrow
\mathbb X
\longrightarrow
\mathbb B
\]
inside the localized homotopy theory of internal categories. The object
\(\mathbb X\) should be read as a derived execution object: it may encode
states, runs, refinements of runs, equivalences between runs, and higher
coherences.

\[
\begin{tikzcd}[column sep=large]
&
\mathbb X
\arrow[dl, "\pi"']
\arrow[dr, "\Omega"]
&
\\
\mathbb A
&
&
\mathbb B
\end{tikzcd}
\]

The map
\[
\pi:\mathbb X\to\mathbb A
\]
records the input semantics of executions, and the map
\[
\Omega:\mathbb X\to\mathbb B
\]
records the observed output behavior. The main result of the chapter is
that, after fixing \(\pi\), the remaining output datum is exactly a point
of the derived mapping space
\[
\operatorname{Map}_{\mathcal C}(\mathbb X,\mathbb B).
\]
This point is what we call a derived \(\mathbb B\)-valued output
\(1\)-cocycle on \(\mathbb X\).

\section{Joyal--Tierney homotopy theory}
\label{sec:derived-mealy-jt-homotopy-theory}

Throughout this chapter, let
\[
\mathcal E
\]
be a Grothendieck topos. Then \(\mathcal E\) has finite limits, so the
bicategory
\[
\mathbf{Span}(\mathcal E)
\]
and the strict Mealy calculus of the previous chapter are still
available.

The additional structure used in this chapter is the Joyal--Tierney
homotopy theory of internal categories in \(\mathcal E\).

\begin{theorem}[Joyal--Tierney]
\label{thm:joyal-tierney-model-structure-derived}
Let \(\mathcal E\) be a Grothendieck topos. The category
\[
\mathbf{Cat}(\mathcal E)
\]
of internal categories in \(\mathcal E\) admits the Joyal--Tierney model
structure. Its weak equivalences are the Joyal--Tierney equivalences, its
cofibrations are the internal functors which are monomorphisms on
objects, and its fibrant objects are the strong stacks.
\end{theorem}

\begin{proof}
This is the Joyal--Tierney model structure on internal categories in a
topos \cite{JoyalTierneyStrongStacks}. We use it as an external
homotopical input throughout this chapter.
\end{proof}

Here a Joyal--Tierney weak equivalence is an internal categorical
equivalence in the sense of Joyal--Tierney: internally fully faithful and
internally essentially surjective, with essential surjectivity interpreted
in the internal-topos sense used in the Joyal--Tierney model structure.

\begin{definition}
\label{def:jt-model-category-notation-derived}
We write
\[
\mathcal M
:=
\mathbf{Cat}(\mathcal E)_{JT}
\]
for the Joyal--Tierney model category of internal categories in
\(\mathcal E\). We write
\[
W_{JT}
\]
for its weak equivalences.

The Dwyer--Kan localization of the relative category
\[
(\mathcal M,W_{JT})
\]
is denoted
\[
\mathcal C
:=
\mathcal M_\infty
:=
L_{DK}(\mathcal M,W_{JT}).
\]
For internal categories \(\mathbb X,\mathbb B\), we write
\[
\mathbf R\!\operatorname{Map}_{JT}(\mathbb X,\mathbb B)
\]
for the corresponding derived mapping space:
\[
\mathbf R\!\operatorname{Map}_{JT}(\mathbb X,\mathbb B)
:=
\operatorname{Map}_{\mathcal C}(\mathbb X,\mathbb B).
\]
\end{definition}

\paragraph{Foundational convention.}
All derived constructions in this chapter are made in the
\(\infty\)-category
\[
\mathcal C=\mathcal M_\infty.
\]
Thus products, pullbacks, slices, mapping spaces, and homotopy fibers are
taken in \(\mathcal C\), unless explicitly stated otherwise.

When a product such as
\[
\mathbb A\times\mathbb B
\]
is written inside \(\mathcal C\), it denotes the product in the
\(\infty\)-category \(\mathcal C\). Point-set products in
\(\mathcal M\) may be used as representatives only insofar as they
present the corresponding derived products in \(\mathcal C\).

Similarly, a point-set diagram
\[
\mathbb A
\xleftarrow{\pi}
\mathbb X
\xrightarrow{\Omega}
\mathbb B
\]
determines a derived correspondence by the induced map
\[
(\pi,\Omega):\mathbb X\to \mathbb A\times\mathbb B
\]
in \(\mathcal C\). Conversely, a derived correspondence is an object of
the slice \(\infty\)-category
\[
\mathcal C/(\mathbb A\times\mathbb B).
\]

Point-set internal categories and internal functors are therefore used as
representatives of objects and morphisms in the localized homotopy
theory, not as the primary invariant objects.

\begin{remark}
\label{rem:dk-not-only-homotopy-category-derived}
The Dwyer--Kan localization is used because we need full derived mapping
spaces, not only the ordinary homotopy category obtained by formally
inverting \(W_{JT}\). Thus all moduli statements below are statements in
the localized homotopy theory, equivalently in the associated
\(\infty\)-category.
\end{remark}

\paragraph{Point-set pullback convention.}
Composition of derived Mealy morphisms is always defined by pullback in
\[
\mathcal C.
\]
Whenever an ordinary point-set pullback in \(\mathcal M\) is asserted to
present such a derived pullback, this is to be read under one of the
following equivalent safeguards:

\begin{enumerate}
\item the cospan has already been replaced by a point-set representative
which computes the corresponding homotopy pullback;
\item the relevant model-categorical hypotheses, such as a right-proper
pullback computation along the displayed fibration, are available;
\item the assertion is made only after applying an explicit fibrant
diagram replacement of the cospan.
\end{enumerate}

This convention is used below only for point-set representatives. It is
not part of the definition of derived Mealy composition, which is
intrinsically the pullback in \(\mathcal C\).

\section{Slices and point-set fibrational representatives}
\label{sec:slices-point-set-fibrational-representatives}

For an object
\[
\mathbb A\in\mathcal C,
\]
let
\[
\mathcal C/\mathbb A
\]
denote the \(\infty\)-categorical slice. Its objects are derived maps
\[
\mathbb X\to\mathbb A.
\]

On the point-set side, for an internal category
\[
\mathbb A\in\mathcal M,
\]
the ordinary slice category
\[
\mathcal M/\mathbb A
\]
has objects internal functors
\[
\mathbb X\to\mathbb A.
\]
A morphism in the slice is declared to be a weak equivalence when its
underlying morphism in \(\mathcal M\) is a Joyal--Tierney weak
equivalence.

\[
\begin{tikzcd}[column sep=large]
\mathbb X
\arrow[rr, "\sim"]
\arrow[dr, "\pi"']
&
&
\mathbb X^{fib}
\arrow[dl, "\pi^{fib}"]
\\
&
\mathbb A
&
\end{tikzcd}
\]

\begin{proposition}[Point-set fibrations present the derived slice]
\label{prop:jt-fibrations-present-derived-slice}
Let
\[
\mathbb A\in\mathcal M.
\]
Let
\[
\mathbf{Fib}_{JT}(\mathbb A)
\subseteq
\mathcal M/\mathbb A
\]
be the full subcategory spanned by the Joyal--Tierney fibrations
\[
\pi:\mathbb X\to\mathbb A.
\]
Equip \(\mathbf{Fib}_{JT}(\mathbb A)\) with the weak equivalences created
in \(\mathcal M\). Then the inclusion
\[
\mathbf{Fib}_{JT}(\mathbb A)
\hookrightarrow
\mathcal M/\mathbb A
\]
induces an equivalence on Dwyer--Kan localizations:
\[
L_{DK}\bigl(\mathbf{Fib}_{JT}(\mathbb A),W_{JT}\bigr)
\simeq
\mathcal C/\mathbb A.
\]
\end{proposition}

\begin{proof}
The slice category
\[
\mathcal M/\mathbb A
\]
carries the usual slice model structure, with weak equivalences and
fibrations created in \(\mathcal M\). Its fibrant objects are precisely
the maps
\[
\mathbb X\to\mathbb A
\]
which are Joyal--Tierney fibrations.

Every object of \(\mathcal M/\mathbb A\) admits a fibrant replacement in
the slice: factor its structure map as a weak equivalence followed by a
Joyal--Tierney fibration,
\[
\mathbb X
\xrightarrow{\sim}
\mathbb X^{fib}
\longrightarrow
\mathbb A.
\]
Thus the full relative subcategory of fibrant objects is homotopically
essentially surjective in the slice. By the standard comparison between
a model category and its full relative subcategory of fibrant objects,
the inclusion induces an equivalence on Dwyer--Kan localizations:
\[
L_{DK}\bigl(\mathbf{Fib}_{JT}(\mathbb A),W_{JT}\bigr)
\simeq
L_{DK}(\mathcal M/\mathbb A)
\simeq
\mathcal C/\mathbb A.
\]
\end{proof}

\begin{definition}[Point-set fibrational representative]
\label{def:point-set-fibrational-representative}
Let
\[
\mathbb A,\mathbb B\in\mathcal M.
\]
A \textbf{Joyal--Tierney fibrational representative} from
\(\mathbb A\) to \(\mathbb B\) is a point-set span of internal categories
\[
\mathbb A
\xleftarrow{\pi}
\mathbb X
\xrightarrow{\Omega}
\mathbb B
\]
such that
\[
\pi:\mathbb X\to\mathbb A
\]
is a Joyal--Tierney fibration.

It represents the derived correspondence determined by the map
\[
(\pi,\Omega):\mathbb X\to\mathbb A\times\mathbb B
\]
in \(\mathcal C\).
\end{definition}

\[
\begin{tikzcd}[column sep=large]
&
\mathbb X
\arrow[dl, "\pi"', two heads]
\arrow[dr, "\Omega"]
\arrow[d, "{(\pi,\Omega)}" description]
&
\\
\mathbb A
&
\mathbb A\times\mathbb B
\arrow[l, "\operatorname{pr}_{\mathbb A}"]
\arrow[r, "\operatorname{pr}_{\mathbb B}"']
&
\mathbb B
\end{tikzcd}
\]

\begin{remark}
\label{rem:fibrational-representatives-not-primary}
Fibrational representatives are useful point-set representatives, but
they are not the primary definition of a derived Mealy morphism. The
primary definition is made in the localized homotopy theory
\[
\mathcal C.
\]
We do not claim here that every object of
\[
\mathcal C/(\mathbb A\times\mathbb B)
\]
admits such a representative without further point-set replacement
choices. The definition singles out those representatives whose input
leg is already fibrational.
\end{remark}

We shall use the following compatibility between the Joyal--Tierney
fibration notion and internal discrete fibrations.

\begin{proposition}[Discrete fibrations are Joyal--Tierney fibrations]
\label{prop:split-discrete-fibrations-are-jt-fibrations-rewritten}
Every internal discrete fibration of internal categories in
\(\mathcal E\) is a fibration in the Joyal--Tierney model structure.
Consequently, every split internal discrete fibration is a
Joyal--Tierney fibration.
\end{proposition}

\begin{proof}
For internal groupoids, Joyal--Tierney prove that discrete fibrations
have the unique right lifting property with respect to cofibration weak
equivalences. Applying the corresponding internal-category form, or
equivalently applying the groupoid-of-isomorphisms test used in the
Joyal--Tierney fibration criterion, shows that an internal discrete
fibration of internal categories is a Joyal--Tierney fibration. The split
case is immediate.
\end{proof}

\begin{lemma}[Point-set pullbacks along fibrational representatives]
\label{lem:point-set-pullback-along-fibration-derived}
Let
\[
\mathbb X\to\mathbb B
\qquad\text{and}\qquad
\mathbb Y\to\mathbb B
\]
be a cospan in \(\mathcal M\). Suppose that the displayed point-set
cospan is homotopy-pullback-ready in the sense of the point-set pullback
convention above; for example, suppose that the leg
\[
\mathbb Y\to\mathbb B
\]
is a Joyal--Tierney fibration and the relevant right-proper or fibrant
diagram replacement hypothesis is available. Then the ordinary pullback
\[
\mathbb X\times_{\mathbb B}\mathbb Y
\]
is a point-set representative of the corresponding pullback in
\[
\mathcal C.
\]
\end{lemma}

\begin{proof}
This is the standard model-categorical computation of homotopy
pullbacks, applied only after the cospan has been arranged to be a
point-set representative of its derived cospan for which the ordinary
pullback computes the derived pullback. Without such a point-set
hypothesis, the definition of the pullback used in this chapter is the
pullback in
\[
\mathcal C.
\]
\end{proof}

\section{Derived Mealy morphisms}
\label{sec:derived-mealy-morphisms-definition}

Let
\[
\mathbb A,\mathbb B\in\mathcal C.
\]
The derived Mealy calculus from \(\mathbb A\) to \(\mathbb B\) is the
homotopy theory of correspondences
\[
\mathbb A
\longleftarrow
\mathbb X
\longrightarrow
\mathbb B
\]
inside \(\mathcal C\).

\begin{definition}[Derived Mealy moduli theory]
\label{def:derived-mealy-moduli-theory-rewritten}
The \textbf{derived Mealy moduli theory} from \(\mathbb A\) to
\(\mathbb B\) is the slice \(\infty\)-category
\[
\mathsf{Mealy}^h(\mathbb A,\mathbb B)
:=
\mathcal C/(\mathbb A\times\mathbb B).
\]

An object of
\[
\mathsf{Mealy}^h(\mathbb A,\mathbb B)
\]
is called a \textbf{derived Mealy morphism}
\[
\mathbb A\rightsquigarrow^h\mathbb B.
\]
\end{definition}

Thus a derived Mealy morphism is represented by a map
\[
\chi:\mathbb X\to\mathbb A\times\mathbb B
\]
in \(\mathcal C\). By the product universal property, such a map is
equivalently a pair
\[
\pi:\mathbb X\to\mathbb A,
\qquad
\Omega:\mathbb X\to\mathbb B
\]
in \(\mathcal C\). We call
\[
\pi:\mathbb X\to\mathbb A
\]
the \textbf{derived input map} and
\[
\Omega:\mathbb X\to\mathbb B
\]
the \textbf{derived output map}.

\[
\begin{tikzcd}[column sep=large]
&
\mathbb X
\arrow[dl, "\pi"']
\arrow[dr, "\Omega"]
\arrow[d, "\chi" description]
&
\\
\mathbb A
&
\mathbb A\times\mathbb B
\arrow[l, "\operatorname{pr}_{\mathbb A}"]
\arrow[r, "\operatorname{pr}_{\mathbb B}"']
&
\mathbb B
\end{tikzcd}
\]

\begin{remark}[Operational reading]
\label{rem:derived-morphism-operational-reading}
The object \(\mathbb X\) is the derived execution object. The map
\(\pi\) says which input process an execution realizes. The map
\(\Omega\) says which output process is observed along that execution.
Unlike the strict state object \(P\), the object \(\mathbb X\) may have
nontrivial higher homotopy: possible executions may themselves have
equivalences and higher coherences.
\end{remark}

\begin{remark}
\label{rem:derived-morphism-moduli-theoretic-rewritten}
The word ``morphism'' is used moduli-theoretically:
\[
\mathsf{Mealy}^h(\mathbb A,\mathbb B)
\]
is the \(\infty\)-category of derived arrows from \(\mathbb A\) to
\(\mathbb B\). Composition is defined by pullback in
\(\mathcal C\), as described in
Section~\ref{sec:composition-derived-mealy-rewritten}.
Equivalently, the assignment
\[
(\mathbb A,\mathbb B)
\longmapsto
\mathcal C/(\mathbb A\times\mathbb B)
\]
is the hom-\(\infty\)-category of the correspondence
\((\infty,2)\)-category
\[
\operatorname{Corr}(\mathcal C),
\]
whose composition is given by pullback.
\end{remark}

The following elementary identification will be used repeatedly.

\begin{proposition}[Correspondences as a slice over the input projection]
\label{prop:mealy-as-iterated-slice}
There is a canonical equivalence of \(\infty\)-categories
\[
\mathcal C/(\mathbb A\times\mathbb B)
\simeq
(\mathcal C/\mathbb A)_{/\operatorname{pr}_{\mathbb A}},
\]
where
\[
\operatorname{pr}_{\mathbb A}:
\mathbb A\times\mathbb B\to\mathbb A
\]
is the product projection in \(\mathcal C\).
\end{proposition}

\begin{proof}
An object of
\[
\mathcal C/(\mathbb A\times\mathbb B)
\]
is a map
\[
\chi:\mathbb X\to\mathbb A\times\mathbb B.
\]
Composing with the product projections gives maps
\[
\pi:\mathbb X\to\mathbb A,
\qquad
\Omega:\mathbb X\to\mathbb B.
\]
Equivalently, \(\chi\) is a map over \(\mathbb A\)
\[
\pi:\mathbb X\to\mathbb A
\longrightarrow
\operatorname{pr}_{\mathbb A}:\mathbb A\times\mathbb B\to\mathbb A.
\]

\[
\begin{tikzcd}[column sep=large]
\mathbb X
\arrow[rr, "\chi"]
\arrow[dr, "\pi"']
&
&
\mathbb A\times\mathbb B
\arrow[dl, "\operatorname{pr}_{\mathbb A}"]
\\
&
\mathbb A
&
\end{tikzcd}
\]

Thus \(\chi\) is precisely an object of the iterated slice
\[
(\mathcal C/\mathbb A)_{/\operatorname{pr}_{\mathbb A}}.
\]
The same argument applies to mapping spaces: a morphism between two maps
to \(\mathbb A\times\mathbb B\) is exactly a morphism between the
corresponding objects over \(\mathbb A\), compatible with the displayed
maps to \(\operatorname{pr}_{\mathbb A}\). Hence the correspondence
defines an equivalence of \(\infty\)-categories.
\end{proof}

\section{The input projection and its fibers}
\label{sec:input-projection-and-fiber-theorem-derived}

The input projection forgets the output coordinate of a derived Mealy
morphism.

\begin{definition}[Input projection]
\label{def:input-projection-derived-rewritten}
The \textbf{input projection} is the functor
\[
\Pi_{\mathbb A,\mathbb B}:
\mathsf{Mealy}^h(\mathbb A,\mathbb B)
=
\mathcal C/(\mathbb A\times\mathbb B)
\longrightarrow
\mathcal C/\mathbb A
\]
defined by postcomposition with
\[
\operatorname{pr}_{\mathbb A}:
\mathbb A\times\mathbb B\to\mathbb A.
\]
Thus a derived correspondence
\[
\chi:\mathbb X\to\mathbb A\times\mathbb B
\]
is sent to its input map
\[
\pi:=\operatorname{pr}_{\mathbb A}\chi:\mathbb X\to\mathbb A.
\]
\end{definition}

\[
\begin{tikzcd}[column sep=large]
\mathcal C/(\mathbb A\times\mathbb B)
\arrow[r, "\Pi_{\mathbb A,\mathbb B}"]
\arrow[d, "\simeq"']
&
\mathcal C/\mathbb A
\\
(\mathcal C/\mathbb A)_{/\operatorname{pr}_{\mathbb A}}
\arrow[ur, "\mathrm{forget}"']
&
\end{tikzcd}
\]

Under the equivalence of Proposition~\ref{prop:mealy-as-iterated-slice},
this is simply the forgetful functor
\[
(\mathcal C/\mathbb A)_{/\operatorname{pr}_{\mathbb A}}
\longrightarrow
\mathcal C/\mathbb A.
\]

\begin{remark}[Automata interpretation of the input projection]
\label{rem:input-projection-automata-interpretation}
The input projection forgets the observation law and remembers only the
execution semantics over the input process. Thus
\[
\Pi_{\mathbb A,\mathbb B}
\]
forgets
\[
\Omega:\mathbb X\to\mathbb B
\]
and retains
\[
\pi:\mathbb X\to\mathbb A.
\]
Operationally, it asks: what is the space, stack, or internal category
of possible executions over the input system, before assigning any
output behavior?
\end{remark}

We now compute the homotopy fiber of
\[
\Pi_{\mathbb A,\mathbb B}
\]
over a fixed input map
\[
\pi:\mathbb X\to\mathbb A.
\]

\begin{lemma}[Fiber of a slice projection]
\label{lem:fiber-of-slice-projection}
Let
\[
\mathcal D
\]
be an \(\infty\)-category, and let
\[
z\in\mathcal D.
\]
Consider the forgetful functor
\[
\mathcal D_{/z}\to\mathcal D.
\]
For any object
\[
d\in\mathcal D,
\]
the homotopy fiber over \(d\), regarded as an \(\infty\)-groupoid, is
naturally equivalent to the mapping space
\[
\operatorname{Map}_{\mathcal D}(d,z).
\]
\end{lemma}

\begin{proof}
The functor
\[
\mathcal D_{/z}\to\mathcal D
\]
is the right fibration classified by the presheaf
\[
d\longmapsto \operatorname{Map}_{\mathcal D}(d,z).
\]
Therefore its fiber over an object \(d\) is the value of this presheaf at
\(d\), namely
\[
\operatorname{Map}_{\mathcal D}(d,z).
\]
Equivalently, an object of the fiber over \(d\) is precisely a morphism
\[
d\to z,
\]
and higher morphisms in the fiber are exactly the higher paths between
such morphisms.
\end{proof}

\begin{theorem}[Fiber theorem, slice form]
\label{thm:fiber-theorem-slice-form-rewritten}
Let
\[
\pi:\mathbb X\to\mathbb A
\]
be an object of the slice \(\infty\)-category
\[
\mathcal C/\mathbb A.
\]
Then there is a natural equivalence
\[
\operatorname{hofib}_{\pi}
\left(
\Pi_{\mathbb A,\mathbb B}
\right)
\simeq
\operatorname{Map}_{\mathcal C/\mathbb A}
\left(
\pi:\mathbb X\to\mathbb A,\,
\operatorname{pr}_{\mathbb A}:\mathbb A\times\mathbb B\to\mathbb A
\right).
\]
\end{theorem}

\begin{proof}
By Proposition~\ref{prop:mealy-as-iterated-slice},
\[
\mathsf{Mealy}^h(\mathbb A,\mathbb B)
\simeq
(\mathcal C/\mathbb A)_{/\operatorname{pr}_{\mathbb A}}.
\]
Under this equivalence, the input projection
\[
\Pi_{\mathbb A,\mathbb B}
\]
is the forgetful functor
\[
(\mathcal C/\mathbb A)_{/\operatorname{pr}_{\mathbb A}}
\longrightarrow
\mathcal C/\mathbb A.
\]
Applying Lemma~\ref{lem:fiber-of-slice-projection} with
\[
\mathcal D=\mathcal C/\mathbb A
\]
and
\[
z=\operatorname{pr}_{\mathbb A}:\mathbb A\times\mathbb B\to\mathbb A
\]
gives the result.
\end{proof}

The slice statement immediately gives the absolute mapping-space form.

\begin{corollary}[Fiber theorem, absolute form]
\label{cor:fiber-theorem-absolute-form-rewritten}
For every input map
\[
\pi:\mathbb X\to\mathbb A
\]
in \(\mathcal C\), there is a natural equivalence
\[
\operatorname{hofib}_{\pi}
\left(
\Pi_{\mathbb A,\mathbb B}
\right)
\simeq
\operatorname{Map}_{\mathcal C}(\mathbb X,\mathbb B)
=
\mathbf R\!\operatorname{Map}_{JT}(\mathbb X,\mathbb B).
\]
\end{corollary}

\begin{proof}
By Theorem~\ref{thm:fiber-theorem-slice-form-rewritten}, the homotopy
fiber is
\[
\operatorname{Map}_{\mathcal C/\mathbb A}
\left(
\pi,\operatorname{pr}_{\mathbb A}
\right).
\]
A map over \(\mathbb A\)
\[
\mathbb X\to\mathbb A\times\mathbb B
\]
from
\[
\pi:\mathbb X\to\mathbb A
\]
to
\[
\operatorname{pr}_{\mathbb A}:\mathbb A\times\mathbb B\to\mathbb A
\]
is equivalently a map
\[
\mathbb X\to\mathbb B
\]
in \(\mathcal C\), since the component to \(\mathbb A\) is already fixed
to be \(\pi\).

Equivalently, using the product universal property in \(\mathcal C\),
\[
\operatorname{Map}_{\mathcal C}
(\mathbb X,\mathbb A\times\mathbb B)
\simeq
\operatorname{Map}_{\mathcal C}(\mathbb X,\mathbb A)
\times
\operatorname{Map}_{\mathcal C}(\mathbb X,\mathbb B),
\]
and the slice mapping space is the homotopy fiber over the point
\[
\pi\in\operatorname{Map}_{\mathcal C}(\mathbb X,\mathbb A)
\]
of the first projection. This homotopy fiber is
\[
\operatorname{Map}_{\mathcal C}(\mathbb X,\mathbb B).
\]
\end{proof}

\[
\begin{tikzcd}[column sep=large]
\operatorname{hofib}_{\pi}(\Pi_{\mathbb A,\mathbb B})
\arrow[r, "\simeq"]
&
\operatorname{Map}_{\mathcal C/\mathbb A}
\left(
\mathbb X\to\mathbb A,\,
\mathbb A\times\mathbb B\to\mathbb A
\right)
\arrow[r, "\simeq"]
&
\operatorname{Map}_{\mathcal C}(\mathbb X,\mathbb B)
\end{tikzcd}
\]

\begin{remark}[Operational meaning of the fiber theorem]
\label{rem:fiber-theorem-meaning-rewritten}
The fiber theorem says that once the input map
\[
\pi:\mathbb X\to\mathbb A
\]
has been fixed, the remaining derived datum is precisely a derived
output map
\[
\mathbb X\to\mathbb B.
\]
Operationally, fixing \(\pi\) fixes the execution semantics over the
input system. The homotopy fiber of
\[
\Pi_{\mathbb A,\mathbb B}
\]
over \(\pi\) is then the space of possible output observation laws on
that execution object. Its points are output laws, its paths are
equivalences between output laws, and its higher paths are higher
coherences between such equivalences.
\end{remark}

\section{\texorpdfstring{Derived output coordinates as $1$--cocycles}{Derived output coordinates as 1-cocycles}}
\label{sec:derived-output-one-cocycles-rewritten}

The fiber theorem identifies the output coordinate of a derived Mealy
morphism with a point of a derived mapping space. This gives the
homotopical notion of output \(1\)-cocycle.

\begin{definition}[Derived output \(1\)-cocycle]
\label{def:derived-output-one-cocycle-rewritten}
Let
\[
\mathbb X,\mathbb B\in\mathcal C.
\]
A \textbf{\(\mathbb B\)-valued derived output \(1\)-cocycle on
\(\mathbb X\)} is a point of the derived mapping space
\[
\operatorname{Map}_{\mathcal C}(\mathbb X,\mathbb B)
=
\mathbf R\!\operatorname{Map}_{JT}(\mathbb X,\mathbb B).
\]
\end{definition}

\begin{remark}[Terminology]
\label{rem:one-cocycle-terminology}
The term ``\(1\)-cocycle'' is used here in the nonabelian categorical
sense. It is not meant to identify
\[
\operatorname{Map}_{\mathcal C}(\mathbb X,\mathbb B)
\]
with a pre-existing abelian cochain complex. A strict map
\[
\mathbb X\to\mathbb B
\]
is specified by object-level and arrow-level data satisfying identity
and composition equations. These are the categorical analogue of
nonabelian \(1\)-cocycle equations. The derived version replaces the
point-set solution space cut out by these strict equations with the
homotopy-invariant derived mapping space
\[
\operatorname{Map}_{\mathcal C}(\mathbb X,\mathbb B).
\]
Thus a derived output \(1\)-cocycle is a homotopy-coherent
\(\mathbb B\)-valued output law.
\end{remark}

\begin{corollary}[Derived outputs are \(1\)-cocycles]
\label{cor:derived-outputs-are-one-cocycles-rewritten}
Let
\[
\chi:\mathbb X\to\mathbb A\times\mathbb B
\]
be a derived Mealy morphism. Write
\[
\pi:=\operatorname{pr}_{\mathbb A}\chi:\mathbb X\to\mathbb A
\]
and
\[
\Omega:=\operatorname{pr}_{\mathbb B}\chi:\mathbb X\to\mathbb B.
\]
Then the output coordinate
\[
\Omega:\mathbb X\to\mathbb B
\]
is a canonical \(\mathbb B\)-valued derived output \(1\)-cocycle on
\(\mathbb X\).
\end{corollary}

\begin{proof}
By definition, a derived output \(1\)-cocycle on \(\mathbb X\) is a point
of
\[
\operatorname{Map}_{\mathcal C}(\mathbb X,\mathbb B).
\]
The second projection of
\[
\chi:\mathbb X\to\mathbb A\times\mathbb B
\]
is precisely such a point:
\[
\Omega:\mathbb X\to\mathbb B.
\]
Equivalently, by Corollary~\ref{cor:fiber-theorem-absolute-form-rewritten},
the homotopy fiber of the input projection over
\[
\pi:\mathbb X\to\mathbb A
\]
is the derived mapping space
\[
\operatorname{Map}_{\mathcal C}(\mathbb X,\mathbb B).
\]
Thus the remaining datum after fixing \(\pi\) is exactly a derived
output \(1\)-cocycle.
\end{proof}

\begin{remark}[Point-set cocycle representatives]
\label{rem:point-set-cocycle-representatives-rewritten}
In a chosen point-set model for a derived mapping space, a point of
\[
\operatorname{Map}_{\mathcal C}(\mathbb X,\mathbb B)
\]
may be represented by a diagram involving weak equivalences and
replacement data. Such a presentation is not part of the definition of a
derived output \(1\)-cocycle in this chapter. The definition is the point
of the derived mapping space itself.
\end{remark}

\section{Strict cocycle defects}
\label{sec:strict-cocycle-defects-rewritten}

The derived mapping space
\[
\mathbf R\!\operatorname{Map}_{JT}(\mathbb X,\mathbb B)
\]
is the homotopy-correct replacement of the point-set solution space of
strict internal functors
\[
\mathbb X\to\mathbb B.
\]
We now isolate the strict equations whose derived replacement is encoded
by this mapping space.

Let
\[
\mathbb X=(X_0,X_1,s_X,t_X,e_X,m_X)
\]
and
\[
\mathbb B=(B_0,B_1,s_B,t_B,e_B,m_B)
\]
be point-set internal categories in \(\mathcal E\).

\begin{definition}[Typed output \(1\)-cochain]
\label{def:typed-output-one-cochain-rewritten}
A \textbf{typed output \(1\)-cochain} from \(\mathbb X\) to \(\mathbb B\)
consists of maps
\[
c_0:X_0\to B_0,
\qquad
c_1:X_1\to B_1
\]
such that
\[
s_Bc_1=c_0s_X,
\qquad
t_Bc_1=c_0t_X.
\]
Thus a typed output \(1\)-cochain is a morphism of the underlying
internal directed graphs.
\end{definition}

The source-target conditions are the commutativity of
\[
\begin{tikzcd}
X_1 \arrow[r, "c_1"] \arrow[d, "s_X"'] &
B_1 \arrow[d, "s_B"]
\\
X_0 \arrow[r, "c_0"'] &
B_0
\end{tikzcd}
\qquad
\begin{tikzcd}
X_1 \arrow[r, "c_1"] \arrow[d, "t_X"'] &
B_1 \arrow[d, "t_B"]
\\
X_0 \arrow[r, "c_0"'] &
B_0.
\end{tikzcd}
\]

\begin{definition}[Unit defect]
\label{def:unit-defect-rewritten}
The \textbf{unit defect} of a typed output \(1\)-cochain
\[
(c_0,c_1):\mathbb X\to\mathbb B
\]
is the comparison between the two maps
\[
X_0\rightrightarrows B_1
\]
given by
\[
c_1e_X
\qquad
\text{and}
\qquad
e_Bc_0.
\]
The unit defect vanishes if
\[
c_1e_X=e_Bc_0.
\]
\end{definition}

The unit-defect diagram is
\[
\begin{tikzcd}
X_0
\arrow[r, "e_X"]
\arrow[d, "c_0"']
&
X_1
\arrow[d, "c_1"]
\\
B_0
\arrow[r, "e_B"']
&
B_1.
\end{tikzcd}
\]

\begin{definition}[Composition defect]
\label{def:composition-defect-rewritten}
Let
\[
X_2:=X_1\times_{t_X,X_0,s_X}X_1
\]
be the object of composable pairs in \(\mathbb X\). The
\textbf{composition defect} of a typed output \(1\)-cochain
\[
(c_0,c_1):\mathbb X\to\mathbb B
\]
is the comparison between the two maps
\[
X_2\rightrightarrows B_1
\]
given by
\[
c_1m_X
\qquad
\text{and}
\qquad
m_B(c_1\times c_1).
\]
Here
\[
c_1\times c_1:
X_1\times_{t_X,X_0,s_X}X_1
\longrightarrow
B_1\times_{t_B,B_0,s_B}B_1
\]
is the induced map on pullbacks, which is well-defined by the
source-target equations for \((c_0,c_1)\). The composition defect
vanishes if
\[
c_1m_X=m_B(c_1\times c_1).
\]
\end{definition}

The composition-defect diagram is
\[
\begin{tikzcd}[column sep=large]
X_1\times_{X_0}X_1
\arrow[r, "c_1\times c_1"]
\arrow[d, "m_X"']
&
B_1\times_{B_0}B_1
\arrow[d, "m_B"]
\\
X_1
\arrow[r, "c_1"']
&
B_1.
\end{tikzcd}
\]

\begin{proposition}[Strict cocycles as vanishing defects]
\label{prop:strict-cocycles-as-vanishing-defects-rewritten}
A typed output \(1\)-cochain
\[
(c_0,c_1):\mathbb X\to\mathbb B
\]
is the underlying cochain of an internal functor
\[
\mathbb X\to\mathbb B
\]
if and only if its unit defect and composition defect vanish.
\end{proposition}

\begin{proof}
The source-target equations are already included in the definition of a
typed output \(1\)-cochain. To be an internal functor, the pair
\[
(c_0,c_1)
\]
must additionally preserve identities and composition. These two
conditions are exactly
\[
c_1e_X=e_Bc_0
\]
and
\[
c_1m_X=m_B(c_1\times c_1).
\]
Thus the typed \(1\)-cochain is an internal functor precisely when the
unit and composition defects vanish.
\end{proof}

\begin{remark}[Derived correction of defects]
\label{rem:derived-correction-defects-rewritten}
The derived output \(1\)-cocycle space
\[
\mathbf R\!\operatorname{Map}_{JT}(\mathbb X,\mathbb B)
\]
may be regarded as the homotopy-invariant replacement of the point-set
solution space cut out by the strict unit and composition equations. It
records the output coordinate in the localized homotopy theory, rather
than as a merely point-set functoriality condition.
\end{remark}

\section{Composition of derived Mealy morphisms}
\label{sec:composition-derived-mealy-rewritten}

Derived Mealy morphisms compose by pullback in the localized homotopy
theory
\[
\mathcal C.
\]

Let
\[
\chi:\mathbb X\to\mathbb A\times\mathbb B
\]
and
\[
\psi:\mathbb Y\to\mathbb B\times\mathbb C
\]
be derived Mealy morphisms. Write their components as
\[
\pi:\mathbb X\to\mathbb A,
\qquad
\Omega:\mathbb X\to\mathbb B,
\]
and
\[
\rho:\mathbb Y\to\mathbb B,
\qquad
\Lambda:\mathbb Y\to\mathbb C.
\]

\[
\begin{tikzcd}[column sep=large]
\mathbb A
&
\mathbb X
\arrow[l, "\pi"']
\arrow[r, "\Omega"]
&
\mathbb B
&
\mathbb Y
\arrow[l, "\rho"']
\arrow[r, "\Lambda"]
&
\mathbb C.
\end{tikzcd}
\]

\begin{definition}[Composite derived Mealy morphism]
\label{def:composite-derived-mealy-rewritten}
The composite
\[
\psi\circ\chi:\mathbb A\rightsquigarrow^h\mathbb C
\]
is the derived correspondence represented by the pullback in
\(\mathcal C\)
\[
\mathbb X\times_{\mathbb B}\mathbb Y
\]
of
\[
\Omega:\mathbb X\to\mathbb B
\qquad
\text{and}
\qquad
\rho:\mathbb Y\to\mathbb B.
\]
Its structure maps are
\[
\mathbb X\times_{\mathbb B}\mathbb Y
\longrightarrow
\mathbb A
\]
given by
\[
\pi\operatorname{pr}_{\mathbb X}
\]
and
\[
\mathbb X\times_{\mathbb B}\mathbb Y
\longrightarrow
\mathbb C
\]
given by
\[
\Lambda\operatorname{pr}_{\mathbb Y}.
\]
Equivalently, the composite is the map
\[
\mathbb X\times_{\mathbb B}\mathbb Y
\longrightarrow
\mathbb A\times\mathbb C
\]
classified by these two components.
\end{definition}

Diagrammatically:
\[
\begin{tikzcd}[column sep=large, row sep=large]
\mathbb X\times_{\mathbb B}\mathbb Y
\arrow[r, "\operatorname{pr}_{\mathbb Y}"]
\arrow[d, "\operatorname{pr}_{\mathbb X}"']
&
\mathbb Y
\arrow[d, "\rho"]
\\
\mathbb X
\arrow[r, "\Omega"']
&
\mathbb B
\end{tikzcd}
\qquad
\begin{tikzcd}[column sep=large, row sep=large]
&
\mathbb X\times_{\mathbb B}\mathbb Y
\arrow[dl, "\pi\operatorname{pr}_{\mathbb X}"']
\arrow[dr, "\Lambda\operatorname{pr}_{\mathbb Y}"]
\arrow[d, dashed, "{(\pi\operatorname{pr}_{\mathbb X},\,\Lambda\operatorname{pr}_{\mathbb Y})}" description]
&
\\
\mathbb A
&
\mathbb A\times\mathbb C
\arrow[l, "\operatorname{pr}_{\mathbb A}"]
\arrow[r, "\operatorname{pr}_{\mathbb C}"']
&
\mathbb C
\end{tikzcd}
\]

\begin{proposition}
\label{prop:composition-derived-mealy-rewritten}
Derived Mealy morphisms compose up to the canonical coherent
equivalences for pullbacks in the \(\infty\)-category
\[
\mathcal C.
\]
The construction is invariant under equivalence of representatives.
\end{proposition}

\begin{proof}
The composite is defined by a pullback in \(\mathcal C\). Pullbacks in an
\(\infty\)-category are invariant under equivalence of diagrams and are
unique up to a contractible space of choices. Therefore replacing either
derived correspondence by an equivalent representative yields an
equivalent pullback object.

Associativity of composition follows from the canonical equivalence
between iterated pullbacks:
\[
(\mathbb X\times_{\mathbb B}\mathbb Y)\times_{\mathbb C}\mathbb Z
\simeq
\mathbb X\times_{\mathbb B}
(\mathbb Y\times_{\mathbb C}\mathbb Z).
\]
The coherence data are the standard coherences for finite limits in an
\(\infty\)-category.
\end{proof}

\begin{remark}[Pipeline interpretation]
\label{rem:derived-pipeline-interpretation-rewritten}
Composition is the derived analogue of pipeline composition. The first
derived Mealy morphism outputs into \(\mathbb B\), while the second is
input-indexed over \(\mathbb B\). The pullback
\[
\mathbb X\times_{\mathbb B}\mathbb Y
\]
is the homotopy-correct execution object of compatible intermediate
behaviors: an execution of the first machine and an execution of the
second machine whose intermediate \(\mathbb B\)-behaviors agree up to
coherent equivalence.
\end{remark}

\section{The stateless subcalculus}
\label{sec:stateless-subcalculus-derived-rewritten}

The previous chapter showed that strict Mealy morphisms carried by
representable spans are precisely internal functors. In the derived
setting, the corresponding notion is that the input map is equivalent to
the identity input map.

\begin{definition}[Stateless derived Mealy morphism]
\label{def:stateless-derived-mealy-rewritten}
A derived Mealy morphism
\[
\chi:\mathbb X\to\mathbb A\times\mathbb B
\]
is \textbf{stateless} if its input map
\[
\pi:\mathbb X\to\mathbb A
\]
is equivalent in
\[
\mathcal C/\mathbb A
\]
to the identity map
\[
1_{\mathbb A}:\mathbb A\to\mathbb A.
\]
\end{definition}

\[
\begin{tikzcd}[column sep=large]
\mathbb X
\arrow[rr, "\simeq"]
\arrow[dr, "\pi"']
&
&
\mathbb A
\arrow[dl, "1_{\mathbb A}"]
\\
&
\mathbb A
&
\end{tikzcd}
\]

\begin{proposition}[Stateless derived Mealy morphisms and derived functors]
\label{prop:stateless-derived-internal-functors-rewritten}
The homotopy fiber of
\[
\Pi_{\mathbb A,\mathbb B}
\]
over the identity input map is the Joyal--Tierney derived mapping theory
of internal functors:
\[
\operatorname{hofib}_{1_{\mathbb A}}
\left(
\Pi_{\mathbb A,\mathbb B}
\right)
\simeq
\operatorname{Map}_{\mathcal C}(\mathbb A,\mathbb B)
=
\mathbf R\!\operatorname{Map}_{JT}(\mathbb A,\mathbb B).
\]
More generally, if an input map
\[
\pi:\mathbb X\to\mathbb A
\]
is equivalent to \(1_{\mathbb A}\) in
\[
\mathcal C/\mathbb A,
\]
then, after choosing such an equivalence, the homotopy fiber over
\(\pi\) is equivalent to
\[
\mathbf R\!\operatorname{Map}_{JT}(\mathbb A,\mathbb B).
\]
\end{proposition}

\begin{proof}
Apply Corollary~\ref{cor:fiber-theorem-absolute-form-rewritten} to the
input map
\[
1_{\mathbb A}:\mathbb A\to\mathbb A.
\]
This gives
\[
\operatorname{hofib}_{1_{\mathbb A}}
\left(
\Pi_{\mathbb A,\mathbb B}
\right)
\simeq
\operatorname{Map}_{\mathcal C}(\mathbb A,\mathbb B).
\]

If
\[
\pi:\mathbb X\to\mathbb A
\]
is equivalent to \(1_{\mathbb A}\) in the slice
\[
\mathcal C/\mathbb A,
\]
then the corresponding base points of
\[
\mathcal C/\mathbb A
\]
lie in the same connected component. A chosen equivalence between them
transports the homotopy fiber over \(\pi\) to the homotopy fiber over
\(1_{\mathbb A}\). Hence the fiber over \(\pi\) is equivalent, after that
choice, to
\[
\operatorname{Map}_{\mathcal C}(\mathbb A,\mathbb B).
\]
\end{proof}

\begin{remark}[Derived representability]
\label{rem:derived-representable-analogue-rewritten}
This is the derived analogue of the representable-carrier theorem from
the previous chapter. Strict representable carriers are replaced by
input maps equivalent to the identity. The remaining output data are
derived internal functors. If the input map is only equivalent to the
identity rather than literally equal to it, the comparison with the
derived functor calculus depends on the chosen equivalence to
\(1_{\mathbb A}\).

Operationally, this is the stateless regime: the execution object has no
hidden state beyond the input object itself, up to derived equivalence.
\end{remark}

\section{Strict Mealy morphisms as split discrete-fibration representatives}
\label{sec:strict-mealy-as-split-discrete-derived-rewritten}

We now show how the strict Mealy morphisms of the previous chapter
determine rigid point-set representatives in the derived calculus.

Let
\[
\mathbb A=(A_0,A_1,s_A,t_A,e_A,m_A)
\]
be an internal category in \(\mathcal E\).

\begin{definition}[Right action of an internal category]
\label{def:right-action-internal-category-rewritten}
A \textbf{right action} of \(\mathbb A\) on an object
\[
p:P\to A_0
\]
of the slice \(\mathcal E/A_0\) consists of a map
\[
\delta:
A_1\times_{t_A,A_0,p}P
\longrightarrow
P
\]
such that
\[
p\delta=s_A\pi_1,
\]
and satisfying the unit and associativity laws
\[
\delta(e_A(px),x)=x,
\]
and, for
\[
a:u\to v,
\qquad
b:v\to w,
\qquad
p(x)=w,
\]
\[
\delta(b\circ a,x)
=
\delta(a,\delta(b,x)).
\]
\end{definition}

The action is contravariant in the input category:
\[
a:u\to v,
\qquad
x\in P_v
\quad\Longrightarrow\quad
\delta(a,x)\in P_u.
\]

\[
\begin{tikzcd}[column sep=large]
P_u
&
P_v
\arrow[l, dashed, "{a^\ast}"']
\\
u
\arrow[r, "a"']
&
v
\end{tikzcd}
\qquad
a^\ast(x)=\delta(a,x).
\]

\begin{definition}[Morphisms of right actions]
\label{def:morphisms-right-actions}
Let
\[
(P,p,\delta)
\qquad\text{and}\qquad
(P',p',\delta')
\]
be right actions of \(\mathbb A\). A morphism of right actions
\[
h:(P,p,\delta)\to(P',p',\delta')
\]
is a map
\[
h:P\to P'
\]
in \(\mathcal E\) such that
\[
p'h=p
\]
and
\[
h(\delta(a,x))=\delta'(a,hx)
\]
in generalized-element notation. Equivalently, the square
\[
\begin{tikzcd}[column sep=large]
A_1\times_{A_0}P
\arrow[r, "{1\times h}"]
\arrow[d, "\delta"']
&
A_1\times_{A_0}P'
\arrow[d, "\delta'"]
\\
P
\arrow[r, "h"']
&
P'
\end{tikzcd}
\]
commutes.
\end{definition}

We write
\[
\mathbf{Act}_r(\mathbb A)
\]
for the resulting category of right actions of \(\mathbb A\).

\subsection{The category of elements}
\label{subsec:category-of-elements-derived-rewritten}

Let
\[
(p:P\to A_0,\delta)
\]
be a right action of \(\mathbb A\). Define an internal category
\[
P\rtimes\mathbb A
\]
as follows.

Its object object is
\[
(P\rtimes\mathbb A)_0=P.
\]
Its arrow object is
\[
(P\rtimes\mathbb A)_1
=
A_1\times_{t_A,A_0,p}P.
\]
A generalized element
\[
(a,x)\in A_1\times_{A_0}P
\]
is read as an arrow
\[
\delta(a,x)\longrightarrow x.
\]
Thus the source and target maps are
\[
s(a,x)=\delta(a,x),
\qquad
t(a,x)=x.
\]
The identity at \(x\in P\) is
\[
(e_A(px),x).
\]

\[
\begin{tikzcd}[column sep=large]
\delta(a,x)
\arrow[r, "{(a,x)}"]
&
x
\\
p(\delta(a,x))=u
\arrow[r, "a"']
&
p(x)=v
\end{tikzcd}
\]

The object of composable pairs in \(P\rtimes\mathbb A\) is canonically
identified with
\[
\left(A_1\times_{t_A,A_0,s_A}A_1\right)
\times_{t_A\pi_2,A_0,p}
P.
\]
A generalized element is a triple
\[
(a,b,x)
\]
where
\[
a:u\to v,
\qquad
b:v\to w,
\qquad
p(x)=w.
\]
Under this identification, the triple \((a,b,x)\) corresponds to the
composable pair
\[
(a,\delta(b,x)),
\qquad
(b,x),
\]
where \((a,\delta(b,x))\) is performed first and \((b,x)\) second.
Composition is defined by
\[
(a,b,x)\longmapsto (b\circ a,x).
\]

\[
\begin{tikzcd}[column sep=large]
\delta(a,\delta(b,x))
\arrow[r, "{(a,\delta(b,x))}"]
&
\delta(b,x)
\arrow[r, "{(b,x)}"]
&
x
\\
u
\arrow[r, "a"']
&
v
\arrow[r, "b"']
&
w
\end{tikzcd}
\]

\begin{proposition}
\label{prop:action-category-of-elements-rewritten}
The data above define an internal category
\[
P\rtimes\mathbb A.
\]
\end{proposition}

\begin{proof}
The source of the composite \((b\circ a,x)\) is
\[
\delta(b\circ a,x),
\]
which equals
\[
\delta(a,\delta(b,x))
\]
by the action associativity law. This is the source of the first arrow in
the corresponding composable pair
\[
(a,\delta(b,x)),
\qquad
(b,x).
\]
The target of the composite is
\[
x,
\]
the target of the second arrow. Thus composition is well typed.

The identity at \(x\) is
\[
(e_A(px),x).
\]
The action unit law gives
\[
\delta(e_A(px),x)=x,
\]
so this arrow has source and target \(x\). The unit laws in
\(\mathbb A\), together with the action unit law, give the left and right
identity axioms in
\[
P\rtimes\mathbb A.
\]

Associativity of composition follows from associativity in
\(\mathbb A\) and the action associativity law. Both ways of composing
three arrows indexed by
\[
a:u\to v,
\qquad
b:v\to w,
\qquad
c:w\to z,
\qquad
p(x)=z
\]
produce the arrow
\[
(c\circ b\circ a,x),
\]
with source determined by repeated application of the action
associativity law. Hence the displayed data define an internal category.
\end{proof}

There is an evident internal functor
\[
\pi_P:P\rtimes\mathbb A\to\mathbb A
\]
given by
\[
(\pi_P)_0=p,
\qquad
(\pi_P)_1(a,x)=a.
\]

\[
\begin{tikzcd}[column sep=large]
A_1\times_{A_0}P
\arrow[r, "{(\pi_P)_1}"]
\arrow[d, "t"']
&
A_1
\arrow[d, "t_A"]
\\
P
\arrow[r, "p"']
&
A_0
\end{tikzcd}
\]

\begin{definition}[Split internal discrete fibration]
\label{def:split-internal-discrete-fibration-rewritten}
An internal functor
\[
F:\mathbb E\to\mathbb A
\]
is an \textbf{internal discrete fibration} if the square
\[
\begin{tikzcd}
E_1 \arrow[r, "F_1"] \arrow[d, "t_E"'] &
A_1 \arrow[d, "t_A"]
\\
E_0 \arrow[r, "F_0"'] &
A_0
\end{tikzcd}
\]
is a pullback.

We call such a discrete fibration \textbf{split} when we regard the
corresponding unique pullback lifts as the chosen cleavage. In the
discrete case this splitting is canonical once the pullback square is
fixed, and its compatibility with identities and composition is strict.
\end{definition}

\begin{remark}
\label{rem:split-discrete-fibration-splitting-canonical-rewritten}
For a discrete fibration, lifts are unique once the pullback square is
fixed. Thus the splitting is canonical in the discrete case. We keep the
word ``split'' to emphasize that the category-of-elements
representatives constructed here carry strictly compatible lifts.
\end{remark}

\begin{proposition}
\label{prop:action-gives-split-discrete-fibration-rewritten}
The functor
\[
\pi_P:P\rtimes\mathbb A\to\mathbb A
\]
is a split internal discrete fibration.
\end{proposition}

\begin{proof}
The arrow object of
\[
P\rtimes\mathbb A
\]
is
\[
A_1\times_{t_A,A_0,p}P.
\]
This is exactly the pullback of
\[
t_A:A_1\to A_0
\]
along
\[
p:P\to A_0.
\]
Therefore the square
\[
\begin{tikzcd}
(P\rtimes\mathbb A)_1 \arrow[r, "(\pi_P)_1"] \arrow[d, "t"'] &
A_1 \arrow[d, "t_A"]
\\
(P\rtimes\mathbb A)_0 \arrow[r, "p"'] &
A_0
\end{tikzcd}
\]
is a pullback.

Explicitly, an arrow
\[
a:u\to v
\]
of \(\mathbb A\) and a state
\[
x\in P
\]
with
\[
p(x)=v
\]
determine a unique lifted arrow
\[
(a,x):\delta(a,x)\to x
\]
over \(a\). The action unit and associativity laws say exactly that
these canonical lifts are strictly compatible with identities and
composition. Hence
\[
\pi_P
\]
is a split internal discrete fibration.
\end{proof}

Conversely, a split internal discrete fibration gives a right action.

\begin{proposition}[Actions and split internal discrete fibrations]
\label{prop:actions-equivalent-split-discrete-fibrations-rewritten}
For fixed
\[
\mathbb A,
\]
the category
\[
\mathbf{Act}_r(\mathbb A)
\]
of right actions of \(\mathbb A\) is equivalent to the category of split
internal discrete fibrations over \(\mathbb A\).
\end{proposition}

\begin{proof}
Given a right action
\[
\delta:A_1\times_{t_A,A_0,p}P\to P,
\]
Proposition~\ref{prop:action-gives-split-discrete-fibration-rewritten}
constructs a split internal discrete fibration
\[
\pi_P:P\rtimes\mathbb A\to\mathbb A.
\]

Conversely, let
\[
F:\mathbb E\to\mathbb A
\]
be a split internal discrete fibration. Put
\[
P:=E_0,
\qquad
p:=F_0:E_0\to A_0.
\]
Since \(F\) is an internal discrete fibration, the square
\[
\begin{tikzcd}
E_1 \arrow[r, "F_1"] \arrow[d, "t_E"'] &
A_1 \arrow[d, "t_A"]
\\
E_0 \arrow[r, "F_0"'] &
A_0
\end{tikzcd}
\]
is a pullback. Hence an arrow of \(\mathbb E\) with codomain \(x\) lying
over an arrow \(a\) of \(\mathbb A\) is uniquely determined by the pair
\[
(a,x)\in A_1\times_{t_A,A_0,p}P.
\]
Define
\[
\delta(a,x)
\]
to be the domain of this lifted arrow. Preservation of identities and
composition by the internal functor \(F\), together with uniqueness of
lifts in the pullback square, gives precisely
\[
\delta(e_A(px),x)=x
\]
and
\[
\delta(b\circ a,x)=\delta(a,\delta(b,x)).
\]
Thus \(p:P\to A_0\) carries a right action of \(\mathbb A\).

A morphism of right actions
\[
h:P\to P'
\]
induces the evident internal functor
\[
P\rtimes\mathbb A\to P'\rtimes\mathbb A
\]
over \(\mathbb A\). Conversely, a functor over \(\mathbb A\) between
split internal discrete fibrations is determined by its map on objects,
and compatibility with arrows is exactly the action-compatibility
condition. Hence the two constructions define inverse equivalences of
categories.
\end{proof}

\subsection{Strict Mealy morphisms}
\label{subsec:strict-mealy-derived-rewritten}

Let
\[
(P,\delta,\omega):\mathbb A\rightsquigarrow\mathbb B
\]
be a strict Mealy morphism in the sense of the previous chapter. Thus
\[
A_0\xleftarrow{p}P\xrightarrow{q}B_0
\]
is a span, and
\[
\delta:
A_1\times_{t_A,A_0,p}P\to P,
\qquad
\omega:
A_1\times_{t_A,A_0,p}P\to B_1
\]
satisfy the strict Mealy typing, unit, and multiplication laws.

The state-update component
\[
\delta
\]
is precisely a right action of \(\mathbb A\) on
\[
p:P\to A_0.
\]
Hence it determines the category of elements
\[
P\rtimes\mathbb A
\]
and a split internal discrete fibration
\[
\pi_P:P\rtimes\mathbb A\to\mathbb A.
\]

The output component determines an internal functor
\[
\Omega_P:P\rtimes\mathbb A\to\mathbb B
\]
by
\[
(\Omega_P)_0=q,
\qquad
(\Omega_P)_1(a,x)=\omega(a,x).
\]

\[
\begin{tikzcd}[column sep=large]
P\rtimes\mathbb A
\arrow[dr, "\pi_P"']
\arrow[rr, "\Omega_P"]
&
&
\mathbb B
\\
&
\mathbb A
&
\end{tikzcd}
\]

\begin{proposition}
\label{prop:omega-output-functor-rewritten}
The Mealy typing, unit, and multiplication laws for
\[
\omega
\]
are precisely the assertion that
\[
\Omega_P:P\rtimes\mathbb A\to\mathbb B
\]
is an internal functor.
\end{proposition}

\begin{proof}
The typing equations for \(\omega\) say
\[
s_B\omega(a,x)=q(\delta(a,x)),
\qquad
t_B\omega(a,x)=q(x).
\]
These are exactly the source-target compatibility equations for the
arrow map of
\[
\Omega_P.
\]

The Mealy unit law
\[
\omega(e_A(px),x)=e_B(qx)
\]
is exactly preservation of identities.

For composable input arrows
\[
a:u\to v,
\qquad
b:v\to w,
\qquad
p(x)=w,
\]
the Mealy multiplication law
\[
\omega(b\circ a,x)
=
\omega(b,x)\circ\omega(a,\delta(b,x))
\]
is exactly preservation of composition in the internal category
\[
P\rtimes\mathbb A.
\]
Hence \(\Omega_P\) is an internal functor, and conversely functoriality
of \(\Omega_P\) gives precisely these equations.
\end{proof}

\begin{theorem}[Strict Mealy morphisms as split discrete-fibration representatives]
\label{thm:strict-mealy-as-split-discrete-fibration-representatives-rewritten}
Every strict Mealy morphism
\[
(P,\delta,\omega):\mathbb A\rightsquigarrow\mathbb B
\]
determines a point-set Joyal--Tierney fibrational representative
\[
\mathbb A
\xleftarrow{\pi_P}
P\rtimes\mathbb A
\xrightarrow{\Omega_P}
\mathbb B
\]
and hence a derived Mealy morphism
\[
\mathbb A\rightsquigarrow^h\mathbb B.
\]

The \(\delta\)-unit and \(\delta\)-multiplication laws are precisely the
strict action, or cleavage, equations making
\[
\pi_P:P\rtimes\mathbb A\to\mathbb A
\]
a split internal discrete fibration. Relative to this input fibration,
the \(\omega\)-unit and \(\omega\)-multiplication laws are precisely the
strict \(\mathbb B\)-valued \(1\)-cocycle equations, equivalently the
functoriality equations for
\[
\Omega_P:P\rtimes\mathbb A\to\mathbb B.
\]
\end{theorem}

\begin{proof}
The state-update component
\[
\delta
\]
is a right action of \(\mathbb A\) on
\[
p:P\to A_0.
\]
By Proposition~\ref{prop:action-category-of-elements-rewritten}, it
determines the internal category
\[
P\rtimes\mathbb A.
\]
By Proposition~\ref{prop:action-gives-split-discrete-fibration-rewritten},
the projection
\[
\pi_P:P\rtimes\mathbb A\to\mathbb A
\]
is a split internal discrete fibration. By
Proposition~\ref{prop:split-discrete-fibrations-are-jt-fibrations-rewritten},
it is a Joyal--Tierney fibration.

By Proposition~\ref{prop:omega-output-functor-rewritten}, the output
component
\[
\omega
\]
determines an internal functor
\[
\Omega_P:P\rtimes\mathbb A\to\mathbb B.
\]
Thus we obtain a point-set Joyal--Tierney fibrational representative
\[
\mathbb A
\xleftarrow{\pi_P}
P\rtimes\mathbb A
\xrightarrow{\Omega_P}
\mathbb B.
\]
This point-set span determines the derived Mealy morphism represented by
the map
\[
(\pi_P,\Omega_P):
P\rtimes\mathbb A
\to
\mathbb A\times\mathbb B
\]
in \(\mathcal C\).

The final assertion is exactly the separation established above:
\(\delta\) supplies the split discrete-fibration structure, while
\(\omega\) supplies the strict output \(1\)-cocycle on the resulting
action category.
\end{proof}

\begin{proposition}[Converse construction]
\label{prop:split-discrete-fibration-output-functor-gives-strict-mealy}
Let
\[
\pi:\mathbb X\to\mathbb A
\]
be a split internal discrete fibration, and let
\[
\Omega:\mathbb X\to\mathbb B
\]
be an internal functor. Then this data determines a strict Mealy morphism
\[
\mathbb A\rightsquigarrow\mathbb B.
\]
\end{proposition}

\begin{proof}
Let
\[
P:=X_0,
\qquad
p:=\pi_0:X_0\to A_0,
\qquad
q:=\Omega_0:X_0\to B_0.
\]
Since \(\pi\) is a split internal discrete fibration, every arrow
\[
a:u\to v
\]
of \(\mathbb A\) and every state
\[
x\in X_0
\]
with
\[
p(x)=v
\]
has a chosen lifted arrow in \(\mathbb X\)
\[
\widetilde a_x:\delta(a,x)\to x
\]
lying over \(a\). Define
\[
\delta(a,x)
\]
to be the domain of this lifted arrow.

The split compatibility of the discrete-fibration lifts with identities
and composition gives
\[
\delta(e_A(px),x)=x
\]
and
\[
\delta(b\circ a,x)
=
\delta(a,\delta(b,x)).
\]
Thus \(\delta\) is a right action of \(\mathbb A\) on
\[
p:P\to A_0.
\]

Define
\[
\omega(a,x):=\Omega_1(\widetilde a_x).
\]
Since \(\Omega\) is an internal functor, the source and target of
\(\omega(a,x)\) are
\[
q(\delta(a,x))
\qquad
\text{and}
\qquad
q(x),
\]
respectively. Preservation of identities gives
\[
\omega(e_A(px),x)=e_B(qx),
\]
and preservation of composition gives
\[
\omega(b\circ a,x)
=
\omega(b,x)\circ\omega(a,\delta(b,x)).
\]
Hence
\[
(P,\delta,\omega)
\]
is a strict Mealy morphism
\[
\mathbb A\rightsquigarrow\mathbb B.
\]
\end{proof}

\begin{remark}[Discrete fibration, not opfibration]
\label{rem:discrete-fibration-not-opfibration-derived-rewritten}
With the orientation used throughout this work, the input projection is
a discrete fibration, not a discrete opfibration. An input arrow
\[
a:u\to v
\]
acts on a state over \(v\) and produces a state over \(u\). Reversing the
orientation of the Mealy structure cell would exchange discrete
fibrations with discrete opfibrations.
\end{remark}

\subsection{Compatibility with strict pipeline composition}
\label{subsec:strict-pipeline-derived-compatibility-rewritten}

The preceding construction is compatible with the strict pipeline
composition from the previous chapter.

Let
\[
(P,\delta_P,\omega_P):\mathbb A\rightsquigarrow\mathbb B
\]
and
\[
(Q,\delta_Q,\omega_Q):\mathbb B\rightsquigarrow\mathbb C
\]
be strict Mealy morphisms, with underlying spans
\[
A_0\xleftarrow{p_P}P\xrightarrow{q_P}B_0,
\qquad
B_0\xleftarrow{p_Q}Q\xrightarrow{q_Q}C_0.
\]
Their strict composite has carrier
\[
P\times_{q_P,B_0,p_Q}Q
\]
and is given by
\[
\delta_{QP}(a,(x,y))
=
\bigl(
\delta_P(a,x),
\delta_Q(\omega_P(a,x),y)
\bigr),
\]
\[
\omega_{QP}(a,(x,y))
=
\omega_Q(\omega_P(a,x),y).
\]

\[
\begin{tikzcd}[column sep=large]
P\rtimes\mathbb A
\arrow[rr, "\Omega_P"]
\arrow[dr, "\pi_P"']
&
&
\mathbb B
&
Q\rtimes\mathbb B
\arrow[l, "\pi_Q"']
\arrow[dl, "\Omega_Q"]
\\
&
\mathbb A
&
\mathbb C
&
\end{tikzcd}
\]

\begin{proposition}[Strict pipeline as pullback of action categories]
\label{prop:strict-pipeline-as-pullback-action-categories-rewritten}
There is a canonical isomorphism of internal categories over
\(\mathbb A\) and \(\mathbb C\):
\[
(P\times_{q_P,B_0,p_Q}Q)\rtimes\mathbb A
\cong
(P\rtimes\mathbb A)
\times_{\mathbb B}
(Q\rtimes\mathbb B),
\]
where the pullback on the right is the ordinary pullback in
\(\mathbf{Cat}(\mathcal E)\) along
\[
\Omega_P:P\rtimes\mathbb A\to\mathbb B
\]
and
\[
\pi_Q:Q\rtimes\mathbb B\to\mathbb B.
\]
Under this identification, the output functor of the strict composite is
the composite projection to
\[
Q\rtimes\mathbb B
\]
followed by
\[
\Omega_Q:Q\rtimes\mathbb B\to\mathbb C.
\]

If this point-set cospan is homotopy-pullback-ready in the sense of
Lemma~\ref{lem:point-set-pullback-along-fibration-derived}---for
example under a right-proper pullback computation along the
Joyal--Tierney fibration
\[
\pi_Q:Q\rtimes\mathbb B\to\mathbb B
\]
or after an explicit fibrant diagram replacement---then this ordinary
pullback presents the corresponding derived pullback in \(\mathcal C\).
\end{proposition}

\begin{proof}
On objects, both sides have object object
\[
P\times_{q_P,B_0,p_Q}Q.
\]
A generalized object is a pair
\[
(x,y)
\]
satisfying
\[
q_P(x)=p_Q(y).
\]

An arrow of
\[
(P\times_{q_P,B_0,p_Q}Q)\rtimes\mathbb A
\]
is a triple
\[
(a,x,y)
\]
where \(a\) is an arrow of \(\mathbb A\) admissible at \(x\), and
\((x,y)\) is a compatible pair of states. Its image in the pullback
category is
\[
\bigl((a,x),(\omega_P(a,x),y)\bigr).
\]
This is well defined because the arrow
\[
\omega_P(a,x)
\]
of \(\mathbb B\) is admissible at \(y\):
\[
p_Q(y)=q_P(x)
\]
and
\[
t_B(\omega_P(a,x))=q_P(x).
\]

\[
\begin{tikzcd}[column sep=large]
\bigl(
\delta_P(a,x),
\delta_Q(\omega_P(a,x),y)
\bigr)
\arrow[r, "{(a,x,y)}"]
&
(x,y)
\\
\delta_P(a,x)
\arrow[r, "{(a,x)}"']
&
x
\\
\delta_Q(\omega_P(a,x),y)
\arrow[r, "{(\omega_P(a,x),y)}"']
&
y
\end{tikzcd}
\]

The source of the arrow
\[
(\omega_P(a,x),y)
\]
in
\[
Q\rtimes\mathbb B
\]
is
\[
\delta_Q(\omega_P(a,x),y).
\]
Therefore the source of the paired arrow is
\[
\left(
\delta_P(a,x),
\delta_Q(\omega_P(a,x),y)
\right),
\]
which is exactly the state-update formula for the strict pipeline
composite.

Conversely, an arrow in the pullback category consists of an arrow
\[
(a,x)
\]
of
\[
P\rtimes\mathbb A
\]
and an arrow
\[
(\beta,y)
\]
of
\[
Q\rtimes\mathbb B
\]
with the same image in \(\mathbb B\). This forces
\[
\beta=\omega_P(a,x),
\]
so the arrow is uniquely of the displayed form.

The source, target, identity, and composition maps are preserved by the
strict Mealy laws for
\[
(P,\delta_P,\omega_P)
\]
and
\[
(Q,\delta_Q,\omega_Q).
\]
For composition, the key equality is the output cocycle law for
\(\omega_P\):
\[
\omega_P(b\circ a,x)
=
\omega_P(b,x)\circ\omega_P(a,\delta_P(b,x)).
\]
This is exactly what makes the second component of the pullback-category
composition agree with the action category of the strict composite.
Hence the displayed correspondence is an isomorphism of internal
categories over
\[
\mathbb A
\]
and
\[
\mathbb C.
\]

Finally,
\[
\pi_Q:Q\rtimes\mathbb B\to\mathbb B
\]
is a split internal discrete fibration, hence a Joyal--Tierney fibration
by Proposition~\ref{prop:split-discrete-fibrations-are-jt-fibrations-rewritten}.
Under the point-set pullback convention of
Lemma~\ref{lem:point-set-pullback-along-fibration-derived}, the ordinary
pullback above then presents the corresponding derived pullback in
\[
\mathcal C.
\]
\end{proof}

\begin{remark}[Strict pipeline composition]
\label{rem:strict-pipeline-derived-composition-rewritten}
Under the assignment of
Theorem~\ref{thm:strict-mealy-as-split-discrete-fibration-representatives-rewritten},
strict pipeline composition is modeled by the ordinary pullback of the
corresponding split discrete-fibration representatives. Passing to the
localized homotopy theory, the derived composite is the pullback in
\[
\mathcal C.
\]
A point-set ordinary pullback may be used as a representative exactly
when the cospan has been arranged to compute that derived pullback.
\end{remark}

\section{Main consolidation}
\label{sec:derived-mealy-main-consolidation-rewritten}

We collect the constructions and identifications of the chapter.

\begin{theorem}[Derived Mealy morphisms and homotopical \(1\)-cocycles]
\label{thm:derived-mealy-main-consolidation-rewritten}
Let \(\mathcal E\) be a Grothendieck topos, and let
\[
\mathcal M=\mathbf{Cat}(\mathcal E)_{JT}
\]
be the Joyal--Tierney model category of internal categories in
\(\mathcal E\). Let
\[
\mathcal C
=
\mathcal M_\infty
=
L_{DK}(\mathcal M,W_{JT})
\]
be its Dwyer--Kan localization.

\begin{enumerate}
\item The derived Mealy moduli theory from
\[
\mathbb A
\]
to
\[
\mathbb B
\]
is the slice \(\infty\)-category
\[
\mathsf{Mealy}^h(\mathbb A,\mathbb B)
=
\mathcal C/(\mathbb A\times\mathbb B).
\]

\item A derived Mealy morphism is equivalently a derived correspondence
\[
\mathbb A
\longleftarrow
\mathbb X
\longrightarrow
\mathbb B
\]
inside \(\mathcal C\).

\item There is a canonical input projection
\[
\Pi_{\mathbb A,\mathbb B}:
\mathsf{Mealy}^h(\mathbb A,\mathbb B)
\longrightarrow
\mathcal C/\mathbb A
\]
given by postcomposition with
\[
\operatorname{pr}_{\mathbb A}:
\mathbb A\times\mathbb B\to\mathbb A.
\]

\item For a fixed input map
\[
\pi:\mathbb X\to\mathbb A,
\]
the homotopy fiber of this projection is the derived mapping space in
the slice:
\[
\operatorname{hofib}_{\pi}
\left(
\Pi_{\mathbb A,\mathbb B}
\right)
\simeq
\operatorname{Map}_{\mathcal C/\mathbb A}
\left(
\pi,\operatorname{pr}_{\mathbb A}
\right).
\]

\item Equivalently,
\[
\operatorname{hofib}_{\pi}
\left(
\Pi_{\mathbb A,\mathbb B}
\right)
\simeq
\operatorname{Map}_{\mathcal C}(\mathbb X,\mathbb B)
=
\mathbf R\!\operatorname{Map}_{JT}(\mathbb X,\mathbb B).
\]

\item Therefore the output coordinate of a derived Mealy morphism is a
canonical point of a derived mapping space. This point is a derived
\(\mathbb B\)-valued output \(1\)-cocycle on \(\mathbb X\).

\item Automata-theoretically, a derived Mealy morphism is a
homotopy-coherent transducer. Its object
\[
\mathbb X
\]
is a derived execution object over the input process
\[
\mathbb A,
\]
and its output coordinate is a homotopy-coherent observation law valued
in
\[
\mathbb B.
\]

\item Strictly, a typed output \(1\)-cochain
\[
(c_0,c_1):\mathbb X\to\mathbb B
\]
is an internal functor if and only if its unit and composition defects
vanish.

\item Derived Mealy morphisms compose, up to coherent equivalence, by
pullback in
\[
\mathcal C.
\]
This is homotopy-coherent pipeline composition: intermediate behaviors
are synchronized by homotopy pullback.

\item The homotopy fiber over the identity input map is the
Joyal--Tierney derived mapping theory of internal functors:
\[
\operatorname{hofib}_{1_{\mathbb A}}
\left(
\Pi_{\mathbb A,\mathbb B}
\right)
\simeq
\operatorname{Map}_{\mathcal C}(\mathbb A,\mathbb B)
=
\mathbf R\!\operatorname{Map}_{JT}(\mathbb A,\mathbb B).
\]
More generally, representatives whose input map is equivalent to
\(1_{\mathbb A}\) have equivalent fibers after choosing such an
equivalence.

\item Strict Mealy morphisms of the previous chapter determine
Joyal--Tierney fibrational point-set representatives
\[
\mathbb A
\xleftarrow{\pi_P}
P\rtimes\mathbb A
\xrightarrow{\Omega_P}
\mathbb B.
\]
The \(\delta\)-laws are the strict action laws giving the split discrete
fibration, and the \(\omega\)-laws are the strict output \(1\)-cocycle
laws for
\[
\Omega_P.
\]

\item Conversely, a split internal discrete fibration over
\[
\mathbb A
\]
together with an output functor to
\[
\mathbb B
\]
recovers a strict Mealy morphism
\[
\mathbb A\rightsquigarrow\mathbb B.
\]

\item For these canonical point-set representatives, strict pipeline
composition is modeled by ordinary pullback of the corresponding split
discrete-fibration representatives. The derived pipeline composite is
the pullback in
\[
\mathcal C.
\]
The ordinary pullback presents that derived pullback whenever the
point-set cospan is homotopy-pullback-ready in the sense fixed above.
\end{enumerate}
\end{theorem}

\begin{proof}
Part (1) is Definition~\ref{def:derived-mealy-moduli-theory-rewritten}.
Part (2) follows from the product universal property in \(\mathcal C\).
Part (3) is Definition~\ref{def:input-projection-derived-rewritten}.
Part (4) is Theorem~\ref{thm:fiber-theorem-slice-form-rewritten}. Part
(5) is Corollary~\ref{cor:fiber-theorem-absolute-form-rewritten}. Part
(6) is Corollary~\ref{cor:derived-outputs-are-one-cocycles-rewritten}.
Part (7) is the interpretation developed throughout the chapter, beginning
with Section~\ref{sec:derived-mealy-orientation}. Part (8) is
Proposition~\ref{prop:strict-cocycles-as-vanishing-defects-rewritten}.
Part (9) is
Proposition~\ref{prop:composition-derived-mealy-rewritten}. Part (10) is
Proposition~\ref{prop:stateless-derived-internal-functors-rewritten}.
Part (11) is
Theorem~\ref{thm:strict-mealy-as-split-discrete-fibration-representatives-rewritten}.
Part (12) is
Proposition~\ref{prop:split-discrete-fibration-output-functor-gives-strict-mealy}.
Part (13) is
Proposition~\ref{prop:strict-pipeline-as-pullback-action-categories-rewritten}.
\end{proof}

\begin{remark}
\label{rem:derived-mealy-conceptual-summary-rewritten}
The derived theory keeps the stateful character of Mealy morphisms while
placing it in the localized homotopy theory of internal categories. The
input coordinate records how executions vary over the input category.
The output coordinate is classified by a derived mapping space, hence by
a derived \(1\)-cocycle.

Thus the hierarchy is
\[
\text{strict Mealy morphisms}
\longrightarrow
\text{split discrete-fibration point-set representatives}
\longrightarrow
\text{derived correspondences in }\mathcal M_\infty.
\]
The first arrow is the canonical assignment constructed in
Theorem~\ref{thm:strict-mealy-as-split-discrete-fibration-representatives-rewritten};
it should not be read as a fully faithful embedding unless the relevant
categories of transformations and their comparison with the localized
derived moduli theory are also specified.

The stateless part of this hierarchy recovers the derived
internal-functor calculus, with the qualification that an input map
merely equivalent to the identity requires a chosen equivalence to
identify its fiber with the fiber over
\[
1_{\mathbb A}.
\]
\end{remark}

\begin{remark}
\label{rem:derived-automata-final-summary}
From the systems-theoretic perspective, the main conceptual point is
that the derived construction does not simply quotient ordinary Mealy
machines by weak equivalence. It enlarges the notion of machine. A
strict Mealy morphism has a rigid state object and a strict transition
law. A derived Mealy morphism has instead a derived execution object
\[
\mathbb X
\]
over the input process
\[
\mathbb A.
\]
Its possible executions, equivalences between executions, and higher
coherences are part of the semantics. The output law is a derived
\(\mathbb B\)-valued \(1\)-cocycle
\[
\mathbb X\to\mathbb B.
\]
Composition of systems is synchronization of execution objects along
intermediate behavior, carried out by homotopy pullback.
\end{remark}